\section{Legal Implications of Resolution Choice}
\label{sec:legal}

The choice of geographic resolution carries legal implications that practitioners must consider.

\paragraph{County-preservation requirements.}
Option~B (tract$\to$county) uses county adjacency as a computational coarsening tool,
but this should not be conflated with legal county-preservation requirements.
Many state constitutions explicitly require minimizing county splits:
California Article~XXI~\S2(d), Colorado Article~V~\S47, Florida Article~III~\S21.
The multi-scale algorithm treats counties as computational units during coarse moves
but does not enforce county-preservation as a hard constraint.
Practitioners in county-preservation states must verify compliance separately using
\texttt{bisect analyze --types splits} after plan generation.

\paragraph{VRA implications of resolution choice.}
The choice of block group (BG) vs.\ census tract resolution can affect minority
population percentages in near-threshold majority-minority districts.
A district that barely satisfies a 50\% minority VAP threshold at the tract level
may fall below threshold when redrawn at the BG level due to finer boundary placement.
Practitioners deploying Option~A (BG$\to$tract) should recompute minority VAP
percentages at the BG level and verify VRA compliance using
\texttt{bisect analyze --types bloc-voting} with BG-level assignments.

\paragraph{Resolution and the uniform distribution.}
Different resolution choices produce plans over different state spaces
(block-group partitions vs.\ tract partitions are distinct mathematical objects).
An ensemble produced at tract resolution and an ensemble produced at BG resolution
are not comparable distributions --- they answer different questions.
Expert witnesses should clearly specify the resolution level in any testimony
about ensemble distributions.

\paragraph{Partisan neutrality of resolution choice.}
Resolution choice is a technical decision but not necessarily a partisan-neutral one.
In states with heavily polarised county structures (e.g., a state where nearly all
counties are strongly Republican or strongly Democratic), Option~B county coarsening
can affect which communities are merged and subsequently which districts lean
which direction.
Resolution choice should be disclosed in any expert report;
practitioners should verify that partisan efficiency gap metrics are stable across
resolution levels before selecting Option~B in contested redistricting proceedings.

\paragraph{Partisan outcome comparison.}
To verify that Option~B multi-scale chains and single-scale tract chains produce
similar plan-space distributions (and thus that resolution choice does not introduce
partisan bias), Table~\ref{tab:tx-partisan} reports mean Democratic seat share and
efficiency gap for the Texas ensembles.

\begin{table}[h]
\centering\small
\caption{TX $k=38$ partisan outcomes: Option~B vs.\ single-scale.
Mean over 8 chains $\times$ 10{,}000 steps = 80{,}000 sampled plans each.}
\label{tab:tx-partisan}
\begin{tabular}{lccc}
\toprule
Method & Mean D-seat share & Efficiency gap & Diff from single-scale \\
\midrule
Single-scale ReCom & 43.2\% & $-2.7\%$ & baseline \\
Multi-scale Option B & 43.0\% & $-2.9\%$ & $-0.2$ pp (not significant) \\
\bottomrule
\end{tabular}
\end{table}

The $0.2$ pp difference in efficiency gap is within sampling noise ($\pm 0.6$ pp at
this ESS level).
We conclude that Option~B does not introduce detectable partisan bias relative to
single-scale ReCom for Texas.

\paragraph{VRA boundary precision.}
Block-group boundaries offer approximately $100$--$200$\,m horizontal precision
in urbanised areas (matching TIGER/Line BG boundary tolerance), compared to
$500$--$1{,}000$\,m for census tract boundaries.
In near-threshold majority-minority districts (minority VAP $45$--$55\%$),
this boundary precision difference can shift estimated VAP by $0.1$--$0.5$
percentage points at district edges, which is legally significant at the
50\% threshold.
Practitioners in covered jurisdictions running near-threshold minority districts
should use block-group resolution to minimise measurement error.

\paragraph{Resolution selection for legal proceedings.}
When multiple resolution levels are technically feasible, practitioners should
select the finest level at which population equality is achievable within the
statutory tolerance.
If the enacted plan was drawn at a specific resolution level, matching that
resolution is the legally defensible default.
Resolution choice must be fixed and documented before ensemble analysis begins;
any change in resolution between analyses should be disclosed in expert reports.
